Man liest die Orte im Abstand von $\tau = \tRtO$ ab und berechnet die mittleren Geschwindigkeiten $\bar v = \Delta s/\tau$ in den einzelnen Sekunden. Bei einer gleichmässig beschleunigten Bewegung wächst die Geschwindigkeit linear mit der Zeit, die mittleren Geschwindigkeiten aufeinanderfolgender Sekunden müssen also jedes Mal um denselben Betrag $a\,\tau$ zunehmen.

\begin{center}
\begin{tabular}{c|cccccc}
$t$ & \tZaO & \tZbO & \tZcO & \tZdO & \tZeO & \tZfO \\
\hline
$s$ & \sRaO & \sRbO & \sRcO & \sRdO & \sReO & \sRfO \\
\hline
$\bar v$ & \multicolumn{2}{c}{\vRaP{0}{3}} & \multicolumn{2}{c}{\vRcP{0}{3}} & \multicolumn{2}{c}{\vReP{0}{3}} \\
& & \multicolumn{2}{c}{\vRbP{0}{3}} & \multicolumn{2}{c}{\vRdP{0}{3}} &
\end{tabular}
\end{center}

Die mittleren Geschwindigkeiten nehmen in jeder Sekunde um denselben Betrag zu (im Rahmen der Ablesegenauigkeit). Die Geschwindigkeit wächst also linear mit der Zeit: Die Bewegung ist $\result{\text{gleichmässig beschleunigt}}$, mit
\begin{align}
 a &= \aReF = \frac{\vRe-\vRa}{4\cdot\tRt} = \aRe \approx \result{\aReP{0}{2}}
\end{align}
Zur Zeit $t = 0$ hat der Körper bereits die Geschwindigkeit $v_0 = \vRnF = \vRn \approx \vRnP{0}{2}$ und befindet sich bei $s_0 = \sRaO$.

Gleichwertig kann man prüfen, dass $\Delta s$ in aufeinanderfolgenden Sekunden immer um denselben Betrag $a\,\tau^2$ zunimmt, oder dass die Kurve eine Parabel $s = s_0+v_0\,t+\frac{1}{2}a\,t^2$ ist.
